% ============================================================================
%   APPENDIX F — DETAILED MODE COLLAPSE ANALYSIS
% ============================================================================

\chapter{Detailed Mode Collapse Analysis}
\label{app:mode-collapse}

This appendix documents in full detail the mode collapse phenomenon 
introduced in Section~\ref{sec:res-mode-collapse}. Four analyses are 
reported: the complete distribution of predictions across the six 
labels for every configuration; the outcomes of the constrained-decoding 
simulation that established the structural nature of the collapse 
(Finding F8); the extended comparison between dominant-class and 
truly-discriminative cases (Finding F9); and a quantitative link 
between prediction diversity and reasoning ability on discriminative 
subsets.

% ----------------------------------------------------------------------------
\section{Complete Prediction Distributions}
\label{sec:app-distributions}

Table~\ref{tab:app-distributions} reports the proportion of predictions 
assigned to each of the six labels for every configuration. The values 
are computed on the full $1{,}015$ patient--criterion pairs. Rows sum to 
$100\%$ up to rounding.

\begin{table}[htbp]
    \centering
    \caption{Complete prediction distributions across the six labels 
        for the twelve configurations. \emph{Incl.}: included. \emph{NIncl.}: 
        not included. \emph{Excl.}: excluded. \emph{NExcl.}: not excluded. 
        \emph{NEI}: not enough information. \emph{NA}: not applicable.}
    \label{tab:app-distributions}
    \small
    \begin{tabular}{@{}lcccccc@{}}
        \toprule
        \textbf{Configuration} & 
        \textbf{Incl.} & \textbf{NIncl.} & 
        \textbf{Excl.} & \textbf{NExcl.} & 
        \textbf{NEI} & \textbf{NA} \\
        \midrule
        Mistral-Base A     &  5.4\% &  1.0\% & 12.5\% & 47.2\% & 33.3\% &  0.6\% \\
        Mistral-Base B     &  9.7\% &  4.9\% &  5.3\% & 45.1\% & 34.2\% &  0.8\% \\
        Mistral-Instruct A &  8.1\% &  3.6\% &  4.6\% & 27.4\% & 47.4\% &  8.9\% \\
        Mistral-Instruct B & 14.8\% &  6.4\% & 32.5\% & 20.3\% & 20.8\% &  5.2\% \\
        Llama-Base A       &  4.3\% &  1.2\% & 45.4\% & 22.9\% & 26.2\% &  0.0\% \\
        Llama-Base B       &  0.8\% &  0.5\% &  0.5\% &  0.6\% & 97.2\% &  0.4\% \\
        Llama-Instruct A   &  0.3\% &  0.4\% &  0.6\% &  0.9\% & 97.8\% &  0.0\% \\
        Llama-Instruct B   &  0.1\% &  0.0\% &  0.1\% &  0.1\% & 99.7\% &  0.0\% \\
        Meditron A         &  6.1\% &  2.1\% & 22.4\% & 64.8\% &  4.6\% &  0.0\% \\
        Meditron B         & 10.4\% &  4.8\% & 21.3\% & 59.3\% &  4.2\% &  0.0\% \\
        BioMistral A       &  0.5\% &  0.4\% &  0.5\% &  0.4\% & 98.2\% &  0.0\% \\
        BioMistral B       &  0.0\% &  0.0\% &  0.0\% &  0.0\% & 100.0\% & 0.0\% \\
        \bottomrule
    \end{tabular}
\end{table}

Five configurations concentrate more than $97\%$ of their predictions 
on the same label (\emph{not enough information}): both BioMistral 
variants, both Llama-Instruct variants, and Llama-Base with Prompt~B. 
Two configurations exhibit partial collapse (Meditron~A and B, with 
\emph{not excluded} accounting for $59.3\%$--$64.8\%$ of predictions). 
The five remaining configurations distribute their predictions more 
evenly across labels, though none exceeds $50\%$ on any single class.

% ----------------------------------------------------------------------------
\section{Constrained-Decoding Simulation (Finding F8)}
\label{sec:app-forced}

To distinguish whether mode collapse reflects a genuine absence of 
clinical reasoning or an artifact of the six-label vocabulary, we 
conducted a controlled simulation. For each collapsed configuration, 
the softmax normalization (Equation~\ref{eq:softmax}) was restricted 
to the two decisive labels of each criterion type: 
$\{$\emph{included}, \emph{not included}$\}$ for inclusion criteria, 
and $\{$\emph{excluded}, \emph{not excluded}$\}$ for exclusion criteria. 
The evasion labels were removed from the candidate set, forcing the 
model to commit to a decisive answer.

Table~\ref{tab:app-forced} reports the prediction distribution under 
this constrained regime, restricted to the five collapsed 
configurations of Section~\ref{sec:app-distributions}. Under 
constrained decoding, all five configurations concentrate their 
predictions on the majority label of the restricted candidate set 
(\emph{not excluded} for exclusion criteria, \emph{included} for 
inclusion criteria), with proportions similar to the original 
collapse.

\begin{table}[htbp]
    \centering
    \caption{Prediction distribution under constrained decoding 
        (evasion labels removed from candidate set). All five 
        collapsed configurations migrate their collapse to the 
        majority label of the restricted set. Accuracy is computed 
        against the expert consensus label.}
    \label{tab:app-forced}
    \small
    \begin{tabular}{@{}lccc@{}}
        \toprule
        \textbf{Configuration} & 
        \textbf{Modal label} & 
        \textbf{Usage rate} & 
        \textbf{Accuracy} \\
        \midrule
        BioMistral A      & not excluded / included & 74.2\% & 0.292 \\
        BioMistral B      & not excluded / included & 76.8\% & 0.294 \\
        Llama-Base B      & not excluded / included & 73.5\% & 0.311 \\
        Llama-Instruct A  & not excluded / included & 78.4\% & 0.298 \\
        Llama-Instruct B  & not excluded / included & 81.1\% & 0.292 \\
        \bottomrule
    \end{tabular}
\end{table}

The constrained-decoding accuracy remains near the base rate of the 
majority class ($47.7\%$ if all exclusion criteria are predicted 
\emph{not excluded}, weighted down by the inclusion cases). The 
collapse thus persists structurally: forcing the model to choose 
between two decisive labels does not restore reasoning ability, it 
simply redirects the concentration of predictions from an evasion 
label to the numerically dominant decisive label. This experiment 
grounds Finding~F8: mode collapse in these configurations is not a 
strategic use of evasion vocabulary but a masking of the absence of 
clinical discrimination.

% ----------------------------------------------------------------------------
\section{Dominant vs.\ Discriminative Cases (Finding F9)}
\label{sec:app-discriminative}

Table~\ref{tab:app-discriminative} extends Table~\ref{tab:res-non-trivial} 
of Section~\ref{sec:res-non-trivial} with the full breakdown across all 
twelve configurations, sorted by discriminative accuracy. The 
dominant-class subset consists of the $777$ pairs whose expert label 
is either \emph{not excluded} or \emph{not enough information}. The 
discriminative subset consists of the $238$ remaining pairs whose 
expert label is one of \emph{included}, \emph{excluded}, 
\emph{not included}, or \emph{not applicable}.

\begin{table}[htbp]
    \centering
    \caption{Extended dominant-vs.-discriminative comparison. The 
        discriminative accuracy on the $238$ pairs is the metric that 
        best reflects clinical reasoning ability. Positive drops 
        indicate that the configuration performs better on the 
        discriminative subset than on the dominant subset.}
    \label{tab:app-discriminative}
    \small
    \begin{tabular}{@{}lccc@{}}
        \toprule
        \textbf{Configuration} & 
        \textbf{Dominant acc.} & 
        \textbf{Discriminative acc.} & 
        \boldmath$\Delta$ \\
        \midrule
        Mistral-Instruct A & 0.183 & 0.719 & $+0.536$ \\
        Mistral-Base A     & 0.462 & 0.613 & $+0.151$ \\
        Mistral-Instruct B & 0.474 & 0.500 & $+0.026$ \\
        Mistral-Base B     & 0.638 & 0.416 & $-0.222$ \\
        Llama-Base A       & 0.219 & 0.265 & $+0.046$ \\
        Llama-Instruct A   & 0.373 & 0.021 & $-0.352$ \\
        Llama-Base B       & 0.398 & 0.013 & $-0.385$ \\
        Llama-Instruct B   & 0.383 & 0.013 & $-0.370$ \\
        BioMistral A       & 0.386 & 0.004 & $-0.382$ \\
        BioMistral B       & 0.377 & 0.000 & $-0.377$ \\
        Meditron A         & 0.846 & 0.000 & $-0.846$ \\
        Meditron B         & 0.822 & 0.000 & $-0.822$ \\
        \bottomrule
    \end{tabular}
\end{table}

Meditron exhibits the widest gap between the two subsets: $84.6\%$ 
accuracy on the $777$ dominant-class pairs but $0.0\%$ on the $238$ 
discriminative pairs. BioMistral shows a similar reversal. These 
patterns are the empirical basis of Finding~F9. Only 
Mistral-Instruct~A reverses the direction and performs substantially 
better on the discriminative subset ($71.9\%$) than on the dominant 
subset ($18.3\%$), suggesting that this configuration is the only 
one whose predictions reflect actual clinical judgment rather than 
majority-class shortcuts.

% ----------------------------------------------------------------------------
\section{Diversity Predicts Discriminative Accuracy}
\label{sec:app-diversity-accuracy}

The prediction diversity metric (normalized Shannon entropy, 
Equation~\ref{eq:diversity}) tracks discriminative accuracy closely. 
Table~\ref{tab:app-diversity-accuracy} reports the two metrics side by 
side, sorted by prediction diversity. The Pearson correlation between 
prediction diversity and discriminative accuracy across the twelve 
configurations is $r = 0.89$ ($p < 0.001$, $n = 12$).

\begin{table}[htbp]
    \centering
    \caption{Prediction diversity and discriminative accuracy, sorted 
        by diversity. Configurations with near-zero diversity produce 
        near-zero accuracy on the $238$ discriminative pairs.}
    \label{tab:app-diversity-accuracy}
    \small
    \begin{tabular}{@{}lcc@{}}
        \toprule
        \textbf{Configuration} & 
        \textbf{Diversity} & 
        \textbf{Discriminative acc.} \\
        \midrule
        Mistral-Instruct B & 0.842 & 0.500 \\
        Mistral-Instruct A & 0.812 & 0.719 \\
        Mistral-Base A     & 0.686 & 0.613 \\
        Mistral-Base B     & 0.577 & 0.416 \\
        Llama-Base A       & 0.525 & 0.265 \\
        Meditron B         & 0.377 & 0.000 \\
        Meditron A         & 0.362 & 0.000 \\
        Llama-Base B       & 0.077 & 0.013 \\
        Llama-Instruct A   & 0.066 & 0.021 \\
        BioMistral A       & 0.055 & 0.004 \\
        Llama-Instruct B   & 0.011 & 0.013 \\
        BioMistral B       & 0.000 & 0.000 \\
        \bottomrule
    \end{tabular}
\end{table}

This correlation supports the design choice of the composite utility 
score (Section~\ref{subsec:met-composite}), in which prediction 
diversity enters multiplicatively. A configuration with high 
calibration but zero diversity cannot produce discriminative 
predictions, and its clinical utility is correctly captured as zero 
by the composite score. Conversely, a configuration with modest 
calibration but high diversity retains the potential for discriminative 
reasoning that can be surfaced through recalibration or ensemble 
methods.

The magnitude of this correlation ($r^2 = 0.79$) indicates that 
prediction diversity alone accounts for nearly $80\%$ of the variance 
in discriminative accuracy across configurations. This near-mechanical 
relationship suggests that mode collapse is not a nuisance property to 
be tolerated but a direct predictor of clinical unsuitability, 
detectable from the prediction distribution alone without any access 
to expert labels. In a practical deployment setting, monitoring the 
entropy of a model's prediction distribution over a moving window of 
recent inputs provides a cheap early-warning signal that does not 
require any ground truth or labelled data.